\chapter{쓰기와 읽기}
% full version: chapters/09_acet.tex

뇌의 기억에는 불편한 특성이 하나 있다. 기억을 저장하는 것이 일을 처리하는 바로 그 연결이라는 점이다. 따로 떨어진 하드 디스크는 없다. 네트워크는 무언가를 떠올릴 때마다 자기 연결을 사용하고, 무언가를 배울 때마다 그 연결을 바꾼다. 새것을 기록하면서 옛것을 망가뜨리지 않으려면 어떻게 해야 할까? 컴퓨터에는 이를 위한 별도의 신호, 곧 “쓰기 허용” 신호가 있다. 뇌에서는 아세틸콜린을 내보내는 방송국 \nt{ACET}이 이 역할을 하는 것으로 보인다.

\section*{두 가지 모드}

그림 몇 장을 저장한 네트워크를 떠올려 보자. 그림의 절반을 보여 주면 나머지를 채워 그린다. 이것이 읽기다. 이번에는 옛 그림과 비슷한 새 그림을 보여 주자. 네트워크가 습관대로 채워 그리면 새것이 아니라 옛것을 보게 되고, 오류를 학습하게 된다. 새것을 기록하려면 네트워크는 잠시 자기 기억을 믿지 말고 세상을 봐야 한다.

\pencilfig{9}{\pencilart{9}}{전선 하나가 네트워크를 전환한다. 새것을 쓸 것인가, 옛것을 떠올릴 것인가.}

아세틸콜린이 하는 일이 정확히 이것이며, 세 가지 효과는 모두 측정되었다. 아세틸콜린은 네트워크가 무언가를 “떠올릴” 때 쓰는 내부 연결을 약하게 하고, 감각 기관에서 들어오는 신호를 강하게 하며, 연결이 쉽게 바뀌도록 한다. 아세틸콜린이 많으면 쓰기 모드다. 네트워크는 세상에 귀를 기울이고 배운다. 적으면 읽기 모드다. 네트워크는 기억에 기대어 나머지를 채워 그린다.

우리 모델에는 둘 다 잘되는 중간 수준이 없다. 쓰려면 내부 연결을 눌러 두어야 하고, 읽으려면 내부 연결이 온전히 필요하다. 그래서 스위치가 꼭 있어야 한다.

\section*{눈을 얼마나 믿을 것인가}

엔지니어들은 오래전에 비슷한 문제를 풀었다. 배가 항로를 따라 갈 때 항해사에게는 배가 있어야 할 위치에 대한 계산값이 있고, 새롭지만 부정확한 관측값이 있다. 각각을 얼마나 믿어야 할까? 최적 필터의 답은 이렇다. 계산을 덜 확신할수록 관측을 더 믿어라. 같은 처방에는 두 번째 내용도 있다. 확신이 적을수록 더 빨리 다시 배워야 한다. 노브 하나가 입력의 가중치와 학습 속도를 함께 조절한다. 아세틸콜린이 하는 일이 바로 이것이다. 여기서 가설이 나온다. 아세틸콜린은 네트워크 자신의 세계 모델이 얼마나 불확실한지를 네트워크에 알린다.

\pencilfig{124}{%
% optimal blending: the less sure the model, the more weight on the observation (and the faster it relearns)
\foreach \dx/\wm/\we/\ttop/\bbot in {0/2.4pt/0.6pt/{모델이 확신함}/{아세틸콜린 적음}, 4.3/0.6pt/2.4pt/{모델이 불확실함}/{아세틸콜린 많음}}{
  \begin{scope}[xshift=\dx cm]
  \draw[penlite=pcViolet, wash=pcViolet, rounded corners=3pt] (0,0.95) rectangle (1.3,1.45); \node[lbl, text=black] at (0.65,1.2) {기억};
  \draw[penlite=pcBlue, wash=pcBlue, rounded corners=3pt] (0,-0.25) rectangle (1.3,0.25); \node[lbl, text=black] at (0.65,0) {눈};
  \draw[dotpen=pcAmber, wash=pcAmber] (2.95,0.6) circle (0.55); \node[lbl, text=black] at (2.95,0.6) {추정};
  \draw[pcViolet, line width=\wm, line cap=round, -{Stealth[length=5pt]}] (1.35,1.2) -- (2.4,0.8);
  \draw[pcBlue, line width=\we, line cap=round, -{Stealth[length=5pt]}] (1.35,0) -- (2.4,0.4);
  \node[font=\small\itshape, text=pcGray, anchor=south] at (1.55,1.6) {\ttop};
  \node[lbl, anchor=north] at (1.55,-0.4) {\bbot};
  \end{scope}}
}{눈을 얼마나 믿을 것인가. 네트워크가 자기 세계 모델을 확신하면 추정은 거의 모두 기억에서 나온다. 확신하지 못하면 관측에서 나오고, 그때는 네트워크가 더 빨리 다시 배운다. 가설에 따르면 이 노브를 돌리는 것이 아세틸콜린이다.}

\section*{분업}

앞 장을 떠올려 보자. 세상이 갑자기 바뀌면, 가설에 따르면 노르아드레날린이 예상 밖의 일을 알아챈다. 아세틸콜린은 새 상황을 다 배울 때까지 네트워크를 쓰기 모드로 붙잡아 둔다. 하나는 “뭔가 바뀌었다”를 알리고, 다른 하나는 “기록하라”를 알린다.

\section*{밤의 재기록}

깊은 잠을 자는 동안에는 아세틸콜린 수준이 낮다. 네트워크는 세상과 끊어진 채 읽기 모드에 있다. 바로 이때 낮에 있었던 일이 네트워크 안에서 “재생”된다. 이것은 측정된 사실이다. 이때 낮에 배운 것이 장기 기억으로 옮겨 쓰인다는 것은 그럴듯한 가설이며, 수면을 다루는 장에서 다시 살펴본다.

\fullversion{9}
